\section{Good labels, the unit group, and rational characters}
\label{sec:goodlabels}

Throughout, \(G\) is a finite group, \(N=\exp(G)\), and
\[
 U_N=(\Z/N)^\times .
\]
We call \(n\geq1\) a \emph{good label} when \(\gcd(n,N)=1\) and a
\emph{ramified label} when every prime divisor of \(n\) divides \(N\).

\subsection{The unit group acts by isometries}

\begin{lemma}
\label{lem:good-labels}
Let \(n\) be a good label and let \(m\) satisfy \(mn\equiv1\pmod N\).  Then:
\begin{enumerate}
\item \(\Psi_n\) depends only on the residue \(n\bmod N\);
\item \(\sigma_n:\Conj(G)\to\Conj(G)\) is a bijection with inverse
\(\sigma_m\), and \(|\sigma_n(C)|=|C|\) for every class \(C\);
\item \(\Psi_n\) is orthogonal for the class form
\eqref{eq:coordinate-class-form}, and
\begin{equation}
 \Psi_n^*=\Psi_n^{-1}=\Psi_m;
 \label{eq:good-adjoint-inverse}
\end{equation}
\item the assignment \(u\mapsto\Psi_u\) is a group homomorphism
\[
 \Psi:U_N\longrightarrow O(V_G,\langle\ ,\ \rangle_G).
\]
\end{enumerate}
\end{lemma}

\begin{proof}
(1)  For \(x\in G\) the order \(|x|\) divides \(N\), so \(x^n\) depends only
on \(n\bmod|x|\), hence only on \(n\bmod N\).

(2)  Since \(x^N=1\) and \(nm\equiv1\pmod N\), we get \((x^n)^m=x^{nm}=x\)
for every \(x\in G\).  Hence \(\sigma_m\sigma_n=\sigma_n\sigma_m=\mathrm{id}\)
and \(\sigma_n\) is a bijection.  For the class sizes, fix a class \(C\) and
consider \(\theta:C\to\sigma_n(C)\), \(\theta(x)=x^n\).  It is well defined
because \(C\) generates a single class of \(n\)th powers.  It is injective:
if \(x^n=y^n\) with \(x,y\in C\) then \(x=(x^n)^m=(y^n)^m=y\).  It is
surjective: if \(z\in\sigma_n(C)\) then \(z^m\in\sigma_m\sigma_n(C)=C\) and
\(\theta(z^m)=z^{mn}=z\).  Hence \(|C|=|\sigma_n(C)|\).

(3)  By \eqref{eq:power-matrix}, \(\Psi_n\) is the permutation operator of
\(\sigma_n\).  Using (2),
\[
 \langle\Psi_nf,\Psi_nh\rangle_G
 =\frac1{|G|}\sum_{C}|C|\,f(\sigma_nC)h(\sigma_nC)
 =\frac1{|G|}\sum_{C'}|\sigma_n^{-1}C'|\,f(C')h(C')
 =\langle f,h\rangle_G,
\]
because \(|\sigma_n^{-1}C'|=|C'|\).  Thus \(\Psi_n\) is orthogonal and
\(\Psi_n^*=\Psi_n^{-1}\).  Finally \(\Psi_m\Psi_n=\Psi_{mn}=\Psi_1=I\) by (1)
and multiplicativity, which gives \eqref{eq:good-adjoint-inverse}.

(4)  Immediate from (1) and (3).
\end{proof}

\begin{proposition}
\label{prop:good-in-commutant}
For every \(u\in U_N\) we have \(\Psi_u\in\Pram(G)'\).  Consequently
\(\Psi(U_N)\) is a finite abelian group of isometries inside
\(\Pram(G)'\).
\end{proposition}

\begin{proof}
Let \(d\mid N\).  Multiplicativity gives
\(\Psi_u\Psi_d=\Psi_{ud}=\Psi_{du}=\Psi_d\Psi_u\).  Taking adjoints of this
identity and using \eqref{eq:good-adjoint-inverse},
\[
 \Psi_d^*\Psi_{u^{-1}}=\Psi_{u^{-1}}\Psi_d^*.
\]
As \(u\mapsto u^{-1}\) is a bijection of \(U_N\), the operator \(\Psi_u\)
commutes with every \(\Psi_d^*\) as well.  By
\cref{prop:minimal-prime-generators} the operators \(\Psi_d,\Psi_d^*\) with
\(d\mid N\) generate \(\Pram(G)\), and an operator commuting with a
generating set commutes with the generated algebra.
\end{proof}

\begin{proposition}
\label{prop:label-splitting}
Every \(n\geq1\) factors uniquely as \(n=n_{\mathrm{r}}n_{\mathrm{g}}\) with
\(n_{\mathrm{r}}\) a ramified label and \(n_{\mathrm{g}}\) a good label, and
then
\begin{equation}
 \Psi_n=\Psi_{n_{\mathrm{r}}}\Psi_{n_{\mathrm{g}}},
 \qquad
 \Psi_{n_{\mathrm{r}}}\in\Pram(G),
 \qquad
 \Psi_{n_{\mathrm{g}}}=\Psi_{u},\ u=n_{\mathrm{g}}\bmod N\in U_N.
 \label{eq:label-splitting}
\end{equation}
Hence
\begin{equation}
 \PAd(G)=\Alg_\Q\bigl(\Pram(G)\cup\Psi(U_N)\bigr).
 \label{eq:full-adams-generated}
\end{equation}
\end{proposition}

\begin{proof}
Let \(n_{\mathrm{r}}\) be the product of the prime powers of \(n\) at primes
dividing \(N\) and \(n_{\mathrm{g}}=n/n_{\mathrm{r}}\); uniqueness is clear.
Multiplicativity gives the first identity, and
\cref{prop:minimal-prime-generators} places \(\Psi_{n_{\mathrm{r}}}\) in
\(\Pram(G)\).  Since \(\Psi(U_N)\) is closed under adjoints by
\eqref{eq:good-adjoint-inverse}, and \(\PAd(G)\) is generated by all
\(\Psi_n\) and their adjoints, \eqref{eq:full-adams-generated} follows.
\end{proof}

\subsection{Rational classes and fields of character values}

Identify \(U_N\) with \(\operatorname{Gal}(\Q(\zeta_N)/\Q)\) by
\(u\mapsto(\zeta_N\mapsto\zeta_N^u)\).  For a rational class
\(R\in\Conj(G)/U_N\) put
\[
 \Gamma_R=\{u\in U_N:\sigma_u(C)=C\ \text{for some, equivalently every, }C\in R\},
 \qquad
 F_R=\Q(\zeta_N)^{\Gamma_R},
\]
the stabilizer being independent of the choice of \(C\) because \(U_N\) is
abelian.  The field \(F_R\) is a classical object: it is the field of
character values on \(R\).  Indeed, for \(g\in C\in R\) and
\(\chi\in\Irr(G)\), the Galois element \(u\) sends \(\chi(g)\) to
\(\chi(g^u)\), so \(u\) fixes every \(\chi(g)\) if and only if
\(\chi(g^u)=\chi(g)\) for all \(\chi\), that is, if and only if
\(g^u\) is conjugate to \(g\)~\cite{Isaacs}.  Hence
\(F_R=\Q(\chi(g):\chi\in\Irr(G))\), and the étale algebra below is the
product of the fields of values of the rational classes.  Set
\begin{equation}
 \mathcal E(G)=\prod_{R\in\Conj(G)/U_N}F_R .
 \label{eq:etale-algebra}
\end{equation}

\begin{theorem}
\label{thm:rational-class-decomposition}
There is an isomorphism of \(\Q[U_N]\)-modules
\begin{equation}
 V_G\;\cong\;\bigoplus_R\Q[U_N/\Gamma_R]\;\cong\;\bigoplus_RF_R=\mathcal E(G),
 \label{eq:etale-decomposition}
\end{equation}
under which \(\Psi_u\) acts as the field automorphism determined by \(u\) on
each factor.  In particular \([F_R:\Q]=|R|\) and
\(\dim_\Q\mathcal E(G)=k(G)\).
\end{theorem}

\begin{proof}
By \cref{lem:good-labels}, \(\Psi_u\) is the permutation operator of
\(\sigma_u\) in the basis \((\delta_C)_C\), so \(V_G\) is the permutation
module \(\Q[\Conj(G)]\) for \(U_N\).  Decomposing \(\Conj(G)\) into orbits
gives \(V_G\cong\bigoplus_R\Q[R]\), and \(\Q[R]\cong\Q[U_N/\Gamma_R]\) because
a transitive action with point stabilizer \(\Gamma_R\) is isomorphic to the
coset action.  Thus \(|R|=[U_N:\Gamma_R]=[F_R:\Q]\).

By the normal basis theorem \(\Q(\zeta_N)\cong\Q[U_N]\) as
\(\Q[U_N]\)-modules.  Taking \(\Gamma_R\)-invariants on both sides and using
that \(U_N\) is abelian, so that \(\Gamma_R\)-invariants of \(\Q[U_N]\) are
spanned by the \(\Gamma_R\)-coset sums, gives
\(F_R=\Q(\zeta_N)^{\Gamma_R}\cong\Q[U_N/\Gamma_R]\) as \(\Q[U_N]\)-modules.
Summing over \(R\) and adding the class counts gives
\(\dim_\Q\mathcal E(G)=\sum_R|R|=k(G)\).
\end{proof}

\begin{corollary}
\label{cor:good-euler-product}
For \(\operatorname{Re}s>1\),
\begin{equation}
 \prod_{p\nmid N}\det\bigl(I-\Psi_pp^{-s}\bigr)^{-1}
 =\prod_R\zeta^{(N)}_{F_R}(s)
 =\zeta^{(N)}_{\mathcal E(G)}(s),
 \label{eq:good-euler-product}
\end{equation}
the superscript denoting omission of the Euler factors at the primes dividing
\(N\).
\end{corollary}

\begin{proof}
Under \eqref{eq:etale-decomposition} the representation of
\(\operatorname{Gal}(\Q(\zeta_N)/\Q)\) on \(V_G\) is
\(\bigoplus_R\operatorname{Ind}_{\Gamma_R}^{U_N}\mathbf 1\), whose Artin
\(L\)-function is \(\prod_R\zeta_{F_R}\) by Artin induction
\cite[\S{}VII]{Washington}.  For \(p\nmid N\) the class of \(p\) in \(U_N\) is
the Frobenius at \(p\), and \(\Psi_p=\Psi_{p\bmod N}\) by
\cref{lem:good-labels}(1), so \(\det(I-\Psi_pp^{-s})^{-1}\) is the Euler
factor at \(p\) of that \(L\)-function.
\end{proof}

\begin{theorem}
\label{thm:dirichlet-factorization}
Let \(\widehat{U_N}=\operatorname{Hom}(U_N,\C^\times)\), and for
\(\chi\in\widehat{U_N}\) put
\begin{equation}
  m_G(\chi) = \#\{R\in\Conj(G)/U_N:\Gamma_R\subseteq\ker\chi\}.
  \label{eq:dirichlet-multiplicity}
\end{equation}
Then the good-label representation has the exact complex decomposition
\begin{equation}
 V_G\otimes_\Q\C
 \cong
 \bigoplus_{\chi\in\widehat{U_N}}\chi^{\oplus m_G(\chi)},
 \label{eq:dirichlet-character-decomposition}
\end{equation}
and its Euler product factors as
\begin{equation}
 \prod_{p\nmid N}\det(I-\Psi_pp^{-s})^{-1}
 =
 \prod_{\chi\in\widehat{U_N}}
 L^{(N)}(s,\chi)^{m_G(\chi)}.
 \label{eq:dirichlet-factorization}
\end{equation}
Here \(\chi\) is viewed as the corresponding Dirichlet character of
conductor dividing \(N\), and \(L^{(N)}\) denotes the Euler product with
the primes dividing \(N\) omitted.  Thus the rational-class stabilizers
determine every Dirichlet factor and its multiplicity without a
Wedderburn computation.
\end{theorem}

\begin{proof}
By \cref{thm:rational-class-decomposition}, the summand attached to \(R\)
is \(\operatorname{Ind}_{\Gamma_R}^{U_N}\mathbf1\).  Since \(U_N\) is
abelian, Frobenius reciprocity gives
\[
 \dim_\C\operatorname{Hom}_{U_N}
 \bigl(\chi,\operatorname{Ind}_{\Gamma_R}^{U_N}\mathbf1\bigr)
 =
 \dim_\C\operatorname{Hom}_{\Gamma_R}
 \bigl(\chi|_{\Gamma_R},\mathbf1\bigr)
 =
 \begin{cases}
 1,&\Gamma_R\subseteq\ker\chi,\\
 0,&\text{otherwise}.
 \end{cases}
\]
Summing over \(R\) proves
\eqref{eq:dirichlet-character-decomposition} and
\eqref{eq:dirichlet-multiplicity}.  For \(p\nmid N\), the operator
\(\Psi_p\) acts on the \(\chi\)-line by \(\chi(p)\); taking determinants
and multiplying the factors \((1-\chi(p)p^{-s})^{-1}\) over the good
primes proves \eqref{eq:dirichlet-factorization}.
\end{proof}

\begin{corollary}
\label{cor:archimedean-factor}
The field \(F_R\) is totally real if and only if \(-1\in\Gamma_R\), that is,
if and only if the classes in \(R\) are real.  Writing \((r_1(R),r_2(R))\) for
the signature of \(F_R\), the completed product
\(\prod_R\zeta_{F_R}\) has archimedean factor
\begin{equation}
 \Gamma_{\mathbf R}(s)^{\sum_R(r_1(R)+r_2(R))}\,
 \Gamma_{\mathbf R}(s+1)^{\sum_Rr_2(R)},
 \qquad
 \Gamma_{\mathbf R}(s)=\pi^{-s/2}\Gamma(s/2),
 \label{eq:archimedean-factor}
\end{equation}
and conductor \(\prod_R|d_{F_R}|\).
\end{corollary}

\begin{proof}
Complex conjugation on \(\Q(\zeta_N)\) is \(\Psi_{-1}\), and \(F_R\) is fixed
pointwise by it if and only if \(-1\in\Gamma_R\); a subfield of a cyclotomic
field is either totally real or totally imaginary, so this is the criterion
for total reality.  By \cref{lem:good-labels}(2) the condition
\(\sigma_{-1}(C)=C\) says that \(C=C^{-1}\).  The displayed factor is
\(\prod_R\Gamma_{\mathbf R}(s)^{r_1(R)}\Gamma_{\mathbf C}(s)^{r_2(R)}\) with
\(\Gamma_{\mathbf C}(s)=\Gamma_{\mathbf R}(s)\Gamma_{\mathbf R}(s+1)\).
\end{proof}

For the Monster, \(k(\M)=194\) with \(172\) rational classes
(\cref{prop:monster-imaginary}), so \(150\) are singletons and \(22\) have two
elements, the latter imaginary quadratic by
\eqref{eq:monster-discriminants}.  Then \(\sum_R(r_1+r_2)=172\),
\(\sum_Rr_2=22\), and \eqref{eq:archimedean-factor} reads
\[
 \Gamma_{\mathbf R}(s)^{172}\,\Gamma_{\mathbf R}(s+1)^{22}.
\]
Here every \(F_R\) has degree at most two.  For \(J_1\) the orbit of size
three at element order \(19\) has \(\Gamma_R\) of order six in
\(U_{19}\cong\Z/18\), and that subgroup contains the involution, so \(F_R\) is
a totally real cubic field contributing \(\Gamma_{\mathbf R}(s)^3\).  Its
block center \(\Q(\zeta_3)\) in \cref{thm:regular-rational-class} is another
field: \(\Q[U_N/\Gamma_R]\cong\Q\oplus W_R\) with \(F_R\) on the left and
\(\Q(\zeta_3)=\End_{\Q[U_N]}(W_R)\) on the right.

\subsection{Ramified labels}

\begin{theorem}
\label{thm:ramified-spectrum}
Let \(p\mid N\), let \(b=v_p(N)\), and let \(\Conj(G)_{p'}\) be the set of
classes of elements of order prime to \(p\).  Then \(\sigma_p\) restricts to a
bijection of \(\Conj(G)_{p'}\), the span
\[
 W_p=\Span_\Q\bigl(\delta_C:C\notin\Conj(G)_{p'}\bigr)
\]
is \(\Psi_p\)-invariant with \(\Psi_p^{\,b}|_{W_p}=0\), and \(\Psi_p\) induces
on \(V_G/W_p\) the permutation operator of \(\sigma_p|_{\Conj(G)_{p'}}\).
Consequently the characteristic polynomial of \(\Psi_p\) is
\[
 X^{\dim W_p}\prod_{\gamma}\bigl(X^{\ell(\gamma)}-1\bigr),
\]
the product being over the cycles \(\gamma\) of
\(\sigma_p|_{\Conj(G)_{p'}}\), and every nonzero eigenvalue of \(\Psi_p\) is a
root of unity.
\end{theorem}

\begin{proof}
If \(|x|=m\) with \(p\nmid m\) then \(|x^p|=m\), so \(\sigma_p\) maps
\(\Conj(G)_{p'}\) into itself; it is injective there, since choosing \(m'\)
with \(pm'\equiv1\bmod m\) recovers \(x\) from \(x^p\) inside
\(\langle x\rangle\), hence bijective.  Conversely if \(p\mid|x|\) then
\(p\mid|x^p|\) unless \(|x|=p\), in which case \(x^p=1\); so for
\(C\notin\Conj(G)_{p'}\) every class in \(\sigma_p^{-1}(C)\) also lies outside
\(\Conj(G)_{p'}\).  By \eqref{eq:power-matrix},
\(\Psi_p\delta_C=\sum_{\sigma_p(C')=C}\delta_{C'}\), so \(W_p\) is
\(\Psi_p\)-invariant.  Moreover \(|x^{p^{b}}|\) is prime to \(p\) for every
\(x\in G\), so \(\sigma_p^{\,b}(\Conj(G))\subseteq\Conj(G)_{p'}\) and
\(\sigma_p^{-b}(C)=\emptyset\) for \(C\notin\Conj(G)_{p'}\); hence
\(\Psi_p^{\,b}|_{W_p}=0\).

The images of the \(\delta_C\) with \(C\in\Conj(G)_{p'}\) form a basis of
\(V_G/W_p\), and \(\Psi_p\) sends \(\delta_C\) to
\(\delta_{\sigma_p^{-1}(C)\cap\Conj(G)_{p'}}\) modulo \(W_p\), which is the
permutation operator of the bijection of the first paragraph.  The
characteristic polynomial is the product of those of \(\Psi_p|_{W_p}\) and of
the induced map, and a permutation of cycle type \((\ell(\gamma))_\gamma\) has
characteristic polynomial \(\prod_\gamma(X^{\ell(\gamma)}-1)\).
\end{proof}

\subsection{Local factors at ramified primes}
\label{sec:ramified-completion}

The good primes give the Euler factors of \(\zeta_{\mathcal E(G)}\) away
from \(N\).  The factors at \(p\mid N\) are determined, as for any
abelian étale algebra, by Frobenius acting on inertia invariants; here both
are built from good-label operators restricted to element-order layers.

For \(m\mid N\), let \(X_m\) be the classes of elements of order \(m\)
and put \(V_m=\operatorname{Map}(X_m,\Q)\).  The action of \(U_N\) on
\(X_m\) factors through \(U_m=(\Z/m)^\times\).  We write
\(\Psi_{u,m}\) for the resulting operator; a lift of \(u\) to \(U_N\)
gives the same operator on \(V_m\).

Fix a prime \(p\).  If \(p\nmid m\), set \(I_{p,m}=1\) and
\(\kappa_{p,m}=p\bmod m\).  If \(m=p^am'\), with \(a>0\) and
\(p\nmid m'\), put
\[
 I_{p,m}=\ker(U_m\longrightarrow U_{m'})
\]
and choose \(\kappa_{p,m}\in U_m\) such that
\begin{equation}
 \kappa_{p,m}\equiv p\pmod{m'},\qquad
 \kappa_{p,m}\equiv1\pmod{p^a}.
 \label{eq:k-p-conditions}
\end{equation}
Let
\begin{equation}
 \Pi_{p,m}=\frac1{|I_{p,m}|}
 \sum_{u\in I_{p,m}}\Psi_{u,m},
 \qquad
 \Phi_{p,m}
 =\Pi_{p,m}\Psi_{\kappa_{p,m},m}\Pi_{p,m}
 \bigm|_{V_m^{I_{p,m}}}.
 \label{eq:frobenius-operator}
\end{equation}
The first operator is the orthogonal projection onto
\(V_m^{I_{p,m}}\).  The second is independent of the choice in
\eqref{eq:k-p-conditions}, since two choices differ by an element of
\(I_{p,m}\).

Define
\begin{equation}
 \mathcal L_p(G,T)=
 \prod_{m\mid N}
 \det\bigl(I-\Phi_{p,m}T\mid V_m^{I_{p,m}}\bigr).
 \label{eq:complete-local-polynomial}
\end{equation}

\begin{theorem}
\label{thm:ramified-completion}
For \(\operatorname{Re}s>1\),
\begin{equation}
 \zeta_{\mathcal E(G)}(s)
 =\prod_R\zeta_{F_R}(s)
 =\prod_p\mathcal L_p(G,p^{-s})^{-1}.
 \label{eq:complete-euler-product}
\end{equation}
If \(p\nmid N\), then
\[
 \mathcal L_p(G,T)=\det(I-\Psi_pT\mid V_G).
\]
If \(p\mid N\), the factors with \(p\nmid m\) are the Frobenius
determinants on the \(p\)-regular order layers, and
\begin{equation}
 C_p(G,T)=
 \prod_{\substack{m\mid N\\p\mid m}}
 \det\bigl(I-\Phi_{p,m}T\mid V_m^{I_{p,m}}\bigr)
 \label{eq:Cp-polynomial}
\end{equation}
is the contribution of the \(p\)-singular layers.
\end{theorem}

\begin{proof}
Every \(U_N\)-orbit \(R\) lies in one \(X_m\).  On
\(\Q[R]\), the local Artin factor of \(F_R\) at \(p\) is
\[
 \det\bigl(I-\operatorname{Frob}_pT
 \mid\Q[R]^{I_{p,m}}\bigr)^{-1}.
\]
For \(p\nmid m\), inertia is trivial and \(p\bmod m\) acts as Frobenius.
For \(p\mid m\), \eqref{eq:k-p-conditions} represents Frobenius modulo
inertia, so \(\Phi_{p,m}\) is its action on the inertia invariants.
Multiplication over the orbits in \(X_m\), then over \(m\mid N\), gives
\(\mathcal L_p(G,T)^{-1}\).  Multiplication over \(p\) gives
\eqref{eq:complete-euler-product}.  If \(p\nmid N\), all inertia groups
are trivial and the direct sum of the \(\Psi_{p,m}\) is \(\Psi_p\).
\end{proof}

\begin{corollary}
\label{cor:complete-euler-product}
The class sizes together with all power maps, that is, the ramified eidetic
object \(\EramObj(G)\), determine every local factor
of \(\zeta_{\mathcal E(G)}\).
\end{corollary}

\begin{corollary}
\label{cor:ramified-degree}
The degree of \(C_p(G,T)\) is
\[
 \delta_p(G)=
 \sum_{\substack{m\mid N\\p\mid m}}
 \#(I_{p,m}\backslash X_m).
\]
If \(r_p(G)\) is the number of \(p\)-singular rational classes and
\(c_p(G)\) the number of \(p\)-singular conjugacy classes, then
\[
 r_p(G)\leq\delta_p(G)\leq c_p(G).
\]
\end{corollary}

\begin{proof}
The fixed space of a permutation group on \(X_m\) has dimension equal to
its number of orbits.  Since \(I_{p,m}\leq U_m\), its orbit partition
refines the rational-class partition and coarsens the partition into
points.
\end{proof}

For the Monster, the ramified degrees are
\[
\begin{array}{c|rrrrrrrrrrrrrrr}
p&2&3&5&7&11&13&17&19&23&29&31&41&47&59&71\\
\hline
\delta_p&
130&90&44&28&11&12&5&4&5&2&3&1&2&1&1
\end{array}
\]
At \(p=2,3\),
\[
 122<130<133,\qquad 88<90<93.
\]

\subsection{Good-label representations}

\begin{theorem}
\label{thm:multiplicity-good-labels}
Let \(A=\Pram(G)\) and \(N=\exp(G)\).  Choose representatives
\(S_1,\ldots,S_t\) for the simple rational \(A\)-modules occurring in
\(V_G\), and write
\[
 D_j=\End_A(S_j),\qquad
 M_j=\operatorname{Hom}_A(S_j,V_G).
\]
Then evaluation gives an \(A\)-module decomposition
\begin{equation}
 V_G\cong\bigoplus_{j=1}^t S_j\otimes_{D_j}M_j,
 \label{eq:double-centralizer-decomposition}
\end{equation}
and restriction to the isotypic components gives
\begin{equation}
 A'\cong\bigoplus_{j=1}^t\End_{D_j}(M_j).
 \label{eq:double-centralizer-commutant}
\end{equation}
For every \(j\), the good labels induce a canonical representation
\begin{equation}
 \rho_j:U_N\longrightarrow
 \End_{D_j}(M_j)^\times,\qquad
 \rho_j(u)(f)=\Psi_u\circ f ,
 \label{eq:good-label-representation}
\end{equation}
with the following properties.
\begin{enumerate}
\item The image is finite and
  \(\rho_j(u)^\dagger=\rho_j(u^{-1})\), where \(\dagger\) is the positive
  involution on \(\End_{D_j}(M_j)\) induced by the class form.
\item Every complex eigenvalue of every \(\rho_j(u)\) is a root of unity,
  and every rational characteristic polynomial of \(\rho_j(u)\) is a
  product of cyclotomic polynomials.
\item The action on the \(j\)-th component is scalar if and only if
  \(\rho_j(U_N)\subseteq Z\bigl(\End_{D_j}(M_j)\bigr)^\times\).
  In that event it defines a character
  \(\lambda_j:U_N\to\mu(K_j)\), where \(K_j=Z(D_j)\).
\item If \(D_j=K_j\) and \(\dim_{K_j}M_j=1\), the action is automatically
  scalar.  Thus the split multiplicity-one theorem below is the rank-one
  specialization of \eqref{eq:good-label-representation}.
\item Put \(a_j=\dim_{D_j}S_j\).  For every \(u\in U_N\),
  \begin{equation}
  \#\{C\in\Conj(G):\sigma_u(C)=C\}
  = \sum_{j=1}^t a_j \operatorname{Tr}_\Q(\rho_j(u)\mid M_j).
  \label{eq:general-lefschetz}
  \end{equation}
\end{enumerate}
\end{theorem}

\begin{proof}
The algebra \(A\) is semisimple by
\cref{thm:rational-bicommutant-structural}.  The rational
double-centralizer theorem therefore gives
\eqref{eq:double-centralizer-decomposition} and
\eqref{eq:double-centralizer-commutant}.  Explicitly, \(S_j\) is a right
\(D_j\)-module through evaluation by \(A\)-endomorphisms, \(M_j\) is the
corresponding left \(D_j\)-multiplicity module, and the evaluation map
\[
 S_j\otimes_{D_j}\operatorname{Hom}_A(S_j,V_G)\longrightarrow V_G,\qquad
 s\otimes f\longmapsto f(s),
\]
identifies the source with the \(S_j\)-isotypic component.

Every \(\Psi_u\) commutes with \(A\) by \cref{prop:good-in-commutant}, so
postcomposition preserves \(\operatorname{Hom}_A(S_j,V_G)\) and is
\(D_j\)-linear.  The equality \(\Psi_{uv}=\Psi_u\Psi_v\) proves that
\(\rho_j\) is a representation.  Because \(U_N\) is finite, its image is
finite.  The identity \(\Psi_u^*=\Psi_{u^{-1}}\) of
\eqref{eq:good-adjoint-inverse} restricts under the double-centralizer
identification to \(\rho_j(u)^\dagger=\rho_j(u^{-1})\), proving (1).

If \(u\) has order \(e\) in \(U_N\), then \(\rho_j(u)^e=1\).  In
characteristic zero the minimal polynomial divides \(x^e-1\), which is
square-free, so \(\rho_j(u)\) is semisimple with roots of unity as
eigenvalues.  Its rational characteristic polynomial is therefore a product
of cyclotomic polynomials.  This proves (2).

For (3),
\[
 Z\bigl(\End_{D_j}(M_j)\bigr)=Z(D_j)=K_j,
\]
and every finite subgroup of \(K_j^\times\) lies in \(\mu(K_j)\).  If
\(D_j=K_j\) and \(\dim_{K_j}M_j=1\), then
\(\End_{K_j}(M_j)=K_j\), proving (4).

Finally, \(\Psi_u\) is the permutation operator of the \(u\)-power action
on conjugacy classes, so its rational trace is the number of fixed classes.
Choose a right \(D_j\)-basis of \(S_j\) of size \(a_j\).  On
\(S_j\otimes_{D_j}M_j\), the operator induced by \(\Psi_u\) is the direct
sum of \(a_j\) copies of \(\rho_j(u)\) as a rational linear operator.
Taking traces in \eqref{eq:double-centralizer-decomposition} gives
\eqref{eq:general-lefschetz}.
\end{proof}

For \(m_j>1\), the good labels define a finite \(\dagger\)-unitary
representation in \(\End_{D_j}(M_j)\).  A scalar character exists if its
image lies in the center.

\subsection{Rational characters on the isotypic components}

Write \(A=\Pram(G)\).  \Cref{prop:good-in-commutant} says that the good
labels live in \(A'\); the structure theory of \cref{sec:structural-theory}
therefore constrains them as soon as we know \(A'\).  The strongest
statement occurs in the multiplicity-free case, which by
\cref{cor:multiplicity-free} is the case \(A'\) commutative.

\begin{theorem}
\label{thm:good-label-characters}
Suppose \(V_G\) is multiplicity free over \(A=\Pram(G)\).  Let
\(E_1,\dots,E_t\) be the primitive idempotents of \(Z(A)\), let
\(A E_k\cong M_{r_k}(D_k)\) be the corresponding simple factors, and let
\(K_k=Z(AE_k)\) be their center fields.  Then \(A'=Z(A)\), and for every
\(u\in U_N\)
\begin{equation}
 \Psi_uE_k=\lambda_k(u)\,E_k,
 \qquad
 \lambda_k:U_N\longrightarrow\mu(K_k),
 \label{eq:good-label-character}
\end{equation}
where \(\mu(K_k)\) is the group of roots of unity of \(K_k\) and
\(\lambda_k\) is a group homomorphism.  Its conductor divides \(N\).  In
particular:
\begin{enumerate}
\item if \(K_k=\Q\) then \(\lambda_k\) is either trivial or a quadratic
Dirichlet character modulo \(N\);
\item for every prime \(p\nmid N\), the compression
\(A_k(p)=\Psi_p|_{V_k}\) satisfies
\begin{equation}
 A_k(p)=\lambda_k(p)\,I_{V_k},
 \qquad
 \det\bigl(I-A_k(p)T\bigr)=\prod_{\iota}\bigl(1-\iota(\lambda_k(p))T\bigr)^{d_k/[K_k:\Q]},
 \label{eq:good-local-factor}
\end{equation}
the product running over the embeddings \(\iota:K_k\hookrightarrow\C\) and
\(d_k=\dim_\Q V_k\).  When \(K_k=\Q\) this reduces to
\(\det(I-A_k(p)T)=(1-\lambda_k(p)T)^{d_k}\).
\end{enumerate}
\end{theorem}

\begin{proof}
By \cref{cor:multiplicity-free} the commutant \(A'\) is commutative and
\(\dim_\Q A'=\dim_\Q Z(A)\).  Since \(Z(A)\subseteq A'\) always, the two
coincide, so \(A'=Z(A)\subseteq A\).

Fix \(u\in U_N\).  By \cref{prop:good-in-commutant}, \(\Psi_u\in A'=Z(A)\).
Multiplying by the primitive central idempotent \(E_k\) gives an element of
\(Z(AE_k)=K_k\), which is \eqref{eq:good-label-character} once we note that
\(K_k\) acts on the ideal \(AE_k\) through the scalar \(\lambda_k(u)\).  The
map \(\lambda_k\) is multiplicative because \(u\mapsto\Psi_u\) is
(\cref{lem:good-labels}(4)) and \(E_k\) is idempotent.  Since \(U_N\) is
finite, \(\lambda_k(u)\) is a root of unity in \(K_k\).  The conductor
divides \(N\) because \(\lambda_k\) is defined on \(U_N\) and
\(\Psi_n\) depends only on \(n\bmod N\) by \cref{lem:good-labels}(1).

(1)  The only roots of unity in \(\Q\) are \(\pm1\), so a homomorphism
\(U_N\to\mu(\Q)=\{\pm1\}\) is trivial or of order two; a homomorphism
\(U_N\to\{\pm1\}\) is a real Dirichlet character modulo \(N\), and
a nontrivial one is quadratic.

(2)  By \eqref{eq:good-label-character}, \(\Psi_p\) acts on
\(V_k=E_kV_G\) as multiplication by the scalar \(\lambda_k(p)\in K_k\).  The
characteristic polynomial of multiplication by \(\alpha\in K_k\) on a
\(K_k\)-vector space of \(\Q\)-dimension \(d_k\) is
\(\prod_\iota(T-\iota(\alpha))^{d_k/[K_k:\Q]}\), which gives
\eqref{eq:good-local-factor}.
\end{proof}

\subsection{Split multiplicity-one blocks}

\begin{theorem}
\label{thm:block-local-character}
Let \(A=\Pram(G)\), let \(E\) be a primitive central idempotent of \(A\) with
center field \(K=Z(AE)\).  Assume that the simple factor is split over its
center, \(AE\cong M_r(K)\), and let \(m\) be the multiplicity of its natural
simple module in \(V_G\).  If \(m=1\), then for every \(u\in U_N\)
\begin{equation}
 \Psi_uE=\lambda_E(u)\,E,
 \qquad
 \lambda_E:U_N\longrightarrow\mu(K),
 \label{eq:block-local-character}
\end{equation}
and \(\lambda_E\) is a homomorphism of conductor dividing \(N\).  No
hypothesis is placed on the other blocks.
\end{theorem}

\begin{proof}
This is the rank-one specialization of
\cref{thm:multiplicity-good-labels}.  Splitness \(AE\cong M_r(K)\) gives
\(D_j=K_j=K\) on the corresponding component, and multiplicity one gives
\(\dim_{K}M_j=1\); by parts (3) and (4) the representation \(\rho_j\) is a
scalar character \(\lambda_E:U_N\to\mu(K)\), which is
\eqref{eq:block-local-character}.  The conductor divides \(N\) because
\(\Psi_n\) depends only on \(n\bmod N\) by \cref{lem:good-labels}(1).
\end{proof}

For a general simple factor \(AE\cong M_r(D)\) with center \(K\), the
double-centralizer theorem gives \(EA'E\cong M_m(D^{\mathrm{op}})\).
For \(m=1\), the good-label action lies in \(D^{\mathrm{op}}\).  If
\(D=K\), it is scalar.

\subsection{The involution on the center}

\begin{lemma}
\label{lem:star-fixes-idempotents}
Every primitive central idempotent \(E\) of \(A=\Pram(G)\) satisfies
\(E^*=E\), and the star restricts to an involution
\(\iota_E\) of the field \(K=Z(AE)\).
\end{lemma}

\begin{proof}
The class form \eqref{eq:coordinate-class-form} is positive definite and
\(A\) is closed under its adjoint, so the isotypic decomposition of \(V_G\)
under \(A\) is orthogonal and \(E\) is the orthogonal projection onto its
isotypic component; an orthogonal projection is self-adjoint.  The star is an
anti-automorphism of \(A\) fixing \(E\), so it preserves the ideal \(AE\) and
hence its center \(K\), on which an anti-automorphism is an automorphism;
it is an involution because the star is.
\end{proof}

\begin{proposition}
\label{prop:star-inverts}
With the hypotheses of \cref{thm:block-local-character},
\begin{equation}
 \lambda_E(u^{-1})=\iota_E\bigl(\lambda_E(u)\bigr)
 =\lambda_E(u)^{-1}
 \qquad(u\in U_N).
 \label{eq:star-inverts}
\end{equation}
Consequently \(\iota_E\) acts on the image group \(\lambda_E(U_N)\) as
inversion.
\end{proposition}

\begin{proof}
Apply the star to \eqref{eq:block-local-character}.  On the left,
\((\Psi_uE)^*=E^*\Psi_u^*=E\Psi_{u^{-1}}=\lambda_E(u^{-1})E\) by
\cref{lem:star-fixes-idempotents} and \eqref{eq:good-adjoint-inverse}.  On
the right, \(\lambda_E(u)\) lies in \(K\subseteq Z(A)\), so
\((\lambda_E(u)E)^*=\iota_E(\lambda_E(u))E\).  Comparing gives the first
equality in \eqref{eq:star-inverts}; the second is multiplicativity of
\(\lambda_E\).
\end{proof}

\begin{corollary}
\label{cor:quadratic-sharp}
Let \(L_E=\Q\bigl(\lambda_E(U_N)\bigr)\subseteq K\) be the field generated by
the realized values.  The following alternatives are mutually exclusive
and exhaustive.
\begin{enumerate}
\item \(\iota_E\) is trivial on \(L_E\).  Then \(\lambda_E^2=\mathbf1\), so
\(\lambda_E\) is trivial or a quadratic Dirichlet character modulo \(N\),
and \(L_E=\Q\).
\item \(\iota_E\) is nontrivial on \(L_E\).  Then \(\lambda_E\) has order
\(e>2\), \(L_E=\Q(\zeta_e)\) is a CM field, and \(\iota_E\) restricts to its
complex conjugation.
\end{enumerate}
In particular \(K=\Q\) implies (1), but (1) is strictly weaker than
\(K=\Q\): what forces quadratic values is triviality of the star on the
generated value field, not rationality of the center.
\end{corollary}

\begin{proof}
By \eqref{eq:star-inverts}, \(\iota_E\) inverts every value.  If \(\iota_E\)
is trivial on \(L_E\) then \(\lambda_E(u)=\lambda_E(u)^{-1}\) for all \(u\),
so every value is \(\pm1\) and \(L_E=\Q\); \cref{thm:good-label-characters}(1)
then identifies \(\lambda_E\).  Otherwise the image is cyclic of some order
\(e\), generated by a primitive \(e\)th root of unity \(\zeta\) with
\(L_E=\Q(\zeta_e)\), and \(\iota_E(\zeta)=\zeta^{-1}\) is a nontrivial
automorphism, so \(e>2\).  A nontrivial automorphism of \(\Q(\zeta_e)\)
inverting \(\zeta_e\) is complex conjugation in every embedding, and a
cyclotomic field of degree greater than one on which complex conjugation is
nontrivial is CM.
\end{proof}

A block carries a character of order greater than two only if its center
admits a complex conjugation, and the values then generate a cyclotomic
subfield of that center.  The center fields occurring in
\cref{sec:characters} are accordingly \(\Q\), \(\Q(\zeta_3)=\Q(\sqrt{-3})\),
and \(\Q(\zeta_5)\); no real quadratic or real quartic center carries a
nontrivial character.

\subsection{Support and character order}

\begin{lemma}
\label{lem:orbit-exponent-bound}
Let \(E\) be a central idempotent of \(A=\Pram(G)\) with \(V_E=EV_G\neq0\)
on which the good labels act by scalars,
\(\Psi_uE=\lambda_E(u)E\).  Let
\[
 S(E)=\{C\in\Conj(G):f(C)\neq0\ \text{for some}\ f\in V_E\}
\]
be the class support of the block.  Then \(S(E)\) is stable under every
\(\sigma_u\), and for each \(u\in U_N\) the multiplicative order of
\(\lambda_E(u)\) divides the order of the permutation
\(\sigma_u|_{S(E)}\).  Consequently the order of \(\lambda_E\) divides the
exponent of the image of \(U_N\) in the symmetric group of \(S(E)\).  In
particular, if every \(\Psi(U_N)\)-orbit meeting \(S(E)\) has at most two
classes, then \(\lambda_E\) is trivial or quadratic.
\end{lemma}

\begin{proof}
For \(f\in V_E\) we have \(\Psi_uf=\lambda_E(u)f\) up to the idempotent:
\(\Psi_u(Ef)=\lambda_E(u)Ef\), and \(\lambda_E(u)\neq0\), so
\(\Psi_uV_E=V_E\).  Since \((\Psi_uf)(C)=f(\sigma_uC)\), the support
satisfies \(\sigma_u^{-1}S(E)\subseteq S(E)\), hence equality by
finiteness.  Now let \(m\) be the order of \(\sigma_u|_{S(E)}\) and take
\(f\in V_E\).  For \(C\in S(E)\),
\((\Psi_u^mf)(C)=f(\sigma_u^mC)=f(C)\); for \(C\notin S(E)\), both sides
vanish because \(S(E)\) is \(\sigma_u\)-stable and \(f\) is supported on
\(S(E)\).  Hence \(\Psi_u^mf=f\), so \(\lambda_E(u)^m=1\).  The exponent
statement follows by taking least common multiples over \(u\), and the
final claim holds because a transitive abelian action is regular: on an
orbit of at most two classes every \(\sigma_u\) squares to the identity,
so \(\lambda_E(u)^2=1\) for every \(u\).
\end{proof}

A character of order greater than two therefore requires a
\(\Psi(U_N)\)-orbit of at least three classes inside the block support.

\begin{proposition}
\label{prop:classical-series}
Let \(n\geq1\).
\begin{enumerate}
\item Every conjugacy class of \(S_n\) is rational.  Hence \(\Psi_u=I\) for
every \(u\in U_N\), every representation \(\rho_j\) of
\cref{thm:multiplicity-good-labels} is trivial, and every good-label
character of \(\Pram(S_n)\) is trivial.
\item Every \(\Psi(U_N)\)-orbit on \(\Conj(A_n)\) has at most two classes.
Hence every \(\rho_j\) has an elementary abelian \(2\)-group as image, and
every scalar block character of \(\Pram(A_n)\) is trivial or quadratic.
\end{enumerate}
\end{proposition}

\begin{proof}
(1)  A permutation and its coprime powers have the same cycle type, hence
are conjugate in \(S_n\).  So \(\sigma_u=\mathrm{id}\) for every good label
\(u\), and \(\Psi_u=I\).

(2)  Let \(x\in A_n\) and let \(u\) be a good label.  Then \(x^u\) has the
same cycle type as \(x\), so the \(A_n\)-class of \(x^u\) lies inside the
\(S_n\)-class of \(x\).  An \(S_n\)-class meets \(A_n\) in at most two
\(A_n\)-classes, so the \(\Psi(U_N)\)-orbit of \([x]\) has at most two
elements.  Every \(\sigma_u\) therefore has order at most two on
\(\Conj(A_n)\), so \(\Psi_u^2=I\), which forces \(\rho_j(u)^2=1\) on every
multiplicity space; the scalar statement is
\cref{lem:orbit-exponent-bound}.
\end{proof}

\subsection{A trace identity}

The following is the all-scalar case of \eqref{eq:general-lefschetz}.

\begin{proposition}
\label{prop:lefschetz}
Let \(E_1,\dots,E_t\) be the primitive central idempotents of
\(A=\Pram(G)\), with center fields \(K_k\) and \(V_k=E_kV_G\), and suppose
each acts by a center scalar, \(\Psi_uE_k=\lambda_k(u)E_k\), for instance
because every simple factor is split over its center and every block has
multiplicity one (\cref{thm:block-local-character}).  Then for every
\(u\in U_N\)
\begin{equation}
 \#\{C\in\Conj(G):\sigma_u(C)=C\}
 =\sum_{k=1}^{t}\frac{d_k}{[K_k:\Q]}\;
 \operatorname{Tr}_{K_k/\Q}\bigl(\lambda_k(u)\bigr),
 \qquad d_k=\dim_\Q V_k .
 \label{eq:lefschetz}
\end{equation}
\end{proposition}

\begin{proof}
By \eqref{eq:power-matrix} the operator \(\Psi_u\) is the permutation matrix
of \(\sigma_u\) in the coordinate basis \((\delta_C)_C\), so its trace is the
number of fixed classes.  Trace is additive over
\(V_G=\bigoplus_kV_k\), and by hypothesis the
restriction \(\Psi_u|_{V_k}\) is multiplication by \(\lambda_k(u)\in K_k\) on
a \(K_k\)-vector space of \(\Q\)-dimension \(d_k\).  The \(\Q\)-trace of
multiplication by \(\alpha\in K_k\) on such a space is
\(\bigl(d_k/[K_k:\Q]\bigr)\operatorname{Tr}_{K_k/\Q}(\alpha)\).
\end{proof}

For a rational center, \cref{thm:good-label-characters}(1) gives a trivial
or quadratic character.  Under the isomorphism
\(U_N\cong\operatorname{Gal}(\Q(\zeta_N)/\Q)\), a quadratic character
\(\lambda_k\) corresponds to a quadratic subfield
\(\Q(\sqrt{D_k})\subseteq\Q(\zeta_N)\), and \(\lambda_k\) is the Kronecker
symbol \(\bigl(\tfrac{D_k}{\cdot}\bigr)\) of the associated fundamental
discriminant \(D_k\), whose prime divisors divide \(N\)~\cite{Davenport}.

\subsection{Saturation for finite groups}

\begin{theorem}
\label{thm:good-label-saturation}
If \(V_G\) is multiplicity free over \(\Pram(G)\), then
\begin{equation}
 \Pram(G)=\PAd(G).
 \label{eq:ram-equals-ad}
\end{equation}
\end{theorem}

\begin{proof}
By \cref{thm:good-label-characters} we have \(A'=Z(A)\) for
\(A=\Pram(G)\), and \(\Psi_u\in A'=Z(A)\subseteq A\) for every \(u\in U_N\).
Hence \(\Psi(U_N)\subseteq\Pram(G)\), and
\eqref{eq:full-adams-generated} gives
\(\PAd(G)\subseteq\Pram(G)\).  The reverse inclusion is trivial.
\end{proof}

\subsection{The rational-class splitting}

Write
\(\Gamma=\Psi(U_N)\subseteq O(V_G)\) and
\[
 V_G^{\Gamma}=\{f\in V_G:\Psi_uf=f\ \text{for all}\ u\in U_N\}.
\]

\begin{proposition}
\label{prop:rational-splitting}
The subspace \(V_G^{\Gamma}\) and its class-orthogonal complement are both
reducing for \(\PAd(G)\), and
\begin{equation}
 \dim_\Q V_G^{\Gamma}=\#\{\text{rational classes of }G\}.
 \label{eq:rational-class-dimension}
\end{equation}
If moreover \(V_G\) is multiplicity free over \(A=\Pram(G)\), then
\begin{equation}
 V_G^{\Gamma}=\bigoplus_{\lambda_k=\mathbf1}V_k,
 \qquad\text{so}\qquad
 \sum_{\lambda_k=\mathbf1}d_k=\#\{\text{rational classes}\},
 \quad
 \sum_{\lambda_k\neq\mathbf1}d_k=k(G)-\#\{\text{rational classes}\}.
 \label{eq:trivial-block-sum}
\end{equation}
\end{proposition}

\begin{proof}
Every element of \(\Gamma\) commutes with \(\Pram(G)\) by
\cref{prop:good-in-commutant} and with \(\Gamma\) itself, so
\(V_G^{\Gamma}\) is invariant under all of \(\Pram(G)\) and all of
\(\Gamma\); by \eqref{eq:full-adams-generated} it is invariant under
\(\PAd(G)\).  Since \(\Gamma\) consists of isometries, the orthogonal
complement of \(V_G^\Gamma\) is \(\Gamma\)-invariant, and it is
\(\PAd(G)\)-invariant because \(\PAd(G)\) is adjoint closed and preserves
\(V_G^\Gamma\).

For \eqref{eq:rational-class-dimension}: \(\Gamma\) acts on \(V_G\) by
permuting the coordinate basis \((\delta_C)_C\) according to the
\(\sigma_u\), so the fixed space has as basis the orbit sums, one for each
\(\Gamma\)-orbit on \(\Conj(G)\).  By \cref{lem:good-labels}(2) and the
remark following it, those orbits are the rational classes.

For \eqref{eq:trivial-block-sum}: assume multiplicity freeness and fix
\(u\in U_N\).  By \eqref{eq:good-label-character}, \(\Psi_u\) acts on
\(V_k\) as multiplication by \(\lambda_k(u)\in K_k\).  If
\(\lambda_k=\mathbf1\) then \(V_k\subseteq V_G^\Gamma\).  If
\(\lambda_k\neq\mathbf1\), choose \(u\) with \(\lambda_k(u)\neq1\); then
\(\lambda_k(u)-1\) is a nonzero element of the field \(K_k\), hence
invertible, so \(\Psi_u-I\) is invertible on the \(K_k\)-module \(V_k\) and
\(V_k\cap V_G^\Gamma=0\).  As \(V_G=\bigoplus_kV_k\) and each \(V_k\) is
\(\Gamma\)-invariant, the fixed space is the sum of the trivial blocks.  The
two displayed sums then follow from \(\sum_kd_k=k(G)\).
\end{proof}

\Cref{prop:rational-splitting} is the general form of the coarse
decomposition that any concrete census exhibits: the total class count
splits as (rational classes) plus (nontrivially twisted dimension), and in
the multiplicity-free case the splitting is by blocks.  For the Monster,
where \(k(G)=194\) and there are \(172\) rational classes, it forces
\(172=\sum_{\lambda_k=\mathbf1}d_k\) and
\(22=\sum_{\lambda_k\neq\mathbf1}d_k\), and \cref{prop:monster-imaginary}
identifies the blocks on either side.

\begin{proposition}
\label{prop:real-rational-parity}
Assume that \(V_G\) is multiplicity free over \(\Pram(G)\).  Then every
nontrivial good-label character \(\lambda_k\) is odd if and only if the
number of real conjugacy classes of \(G\) equals the number of rational
classes of \(G\).  Namely,
\begin{equation}
 \#\{\text{real classes}\}-\#\{\text{rational classes}\}
 =
 \sum_{\substack{\lambda_k\neq\mathbf1\\ \lambda_k(-1)=1}}d_k .
 \label{eq:real-rational-parity}
\end{equation}
\end{proposition}

\begin{proof}
The real classes are the fixed points of inversion on
\(\Conj(G)\), so their number is the dimension of the \(+1\)-eigenspace of
\(\Psi_{-1}\).  By multiplicity freeness, \(\Psi_{-1}\) acts on \(V_k\) as
the scalar \(\lambda_k(-1)\), and hence
\[
 \#\{\text{real classes}\}
 =\sum_{\lambda_k(-1)=1}d_k .
\]
By \eqref{eq:trivial-block-sum}, the sum over the trivial characters is the
number of rational classes.  Subtracting gives
\eqref{eq:real-rational-parity}.  The right-hand side vanishes if and only
if every nontrivial character is odd.
\end{proof}

\subsection{Compressions}

Channel compressions obey the multiplicative law needed to define local
block data without an ordering convention.

\begin{proposition}
\label{prop:compression-multiplicative}
Let \(E\) be a central idempotent of \(A=\Pram(G)\), let \(V_E=EV_G\), and
for \(n\geq1\) put \(A_E(n)=E\Psi_nE|_{V_E}\).  Then \(E\) commutes with
\(\Psi_n\) for every \(n\geq1\), and
\begin{equation}
 A_E(m)A_E(n)=A_E(mn)=A_E(n)A_E(m)
 \qquad(m,n\geq1).
 \label{eq:compression-multiplicative}
\end{equation}
\end{proposition}

\begin{proof}
For a ramified label \(n\) we have \(\Psi_n\in A\) and \(E\in Z(A)\), so
\(E\Psi_n=\Psi_nE\).  For a good label \(n\) we have \(\Psi_n\in A'\) by
\cref{prop:good-in-commutant} and \(E\in A\), so again
\(E\Psi_n=\Psi_nE\).  A general \(n\) is a product of the two kinds by
\cref{prop:label-splitting}.  Therefore
\[
 E\Psi_mE\cdot E\Psi_nE=E^3\Psi_m\Psi_n=E\Psi_{mn},
\]
and restricting to \(V_E\) gives \eqref{eq:compression-multiplicative}; the
symmetry in \(m,n\) holds because \(\Psi_m\Psi_n=\Psi_{mn}=\Psi_n\Psi_m\).
\end{proof}
